% Separate continuation module: do not silently identify its Psi with a
% source's rho having a different monodromy convention.
\subsection*{A single mixed-component duplication defect}

\paragraph{MD.1 (Proved): conventions and the peripheral identity.}
In this module $\Phi=(1,f)$ acts by
$x\mapsto x$, $y\mapsto f^{-1}yf$, and
$[x,y]=xyx^{-1}y^{-1}$. We write
\[
 f\equiv A_f(U,V)[x,y]\pmod{F''},\qquad
 B_f(U,V)=1+(U-1)(V-1)A_f(U,V).
\]
The module action is left conjugation. We denote by $\Psi_n(\Phi)$
the coefficient $m_0$ in KM.1, without assigning it an arithmetic
sign by the symbol $\rho_n$. The conventions needed below are
\[
\begin{array}{c|l}
\text{symbol}&\text{definition}\\\hline
\Psi_n&m_0\text{ in }\ker(F\to C_n)^{\rm ab}\text{ with }y_0=y,\\
\rho&\beta_\Phi^{\rm ab}=(-\rho,-\rho),\\
\kappa_d&\alpha_d(\Phi)^{\rm ab}=\kappa_d e_5\quad(\Phi\in\mathcal H_5),\\
k_n&\sigma(n^{1/m})/n^{1/m}=\zeta_m^{k_n(\sigma)}.
\end{array}
\]
Here $\mathcal H_5=\widehat{GT}\cap\mathcal H$.
Ihara II, Theorem 1, writes the arithmetic identity with $n^{-r}$,
so his $\tau_n^{(0)}$ is $-k_n$ in that monodromy convention.
Furusho (2003), p.100, explicitly records the corresponding minus
sign for the NT characters. The older NS (2000) formulation uses
its frame-defined $\rho_2$ and relates it to its own Alexander
coefficient in Proposition 5.3, on its specified subgroup.
These source symbols are not identified across markings here.

The simultaneous substitution $x\mapsto x^{-1}$,
$y\mapsto y^{-1}$ sends $y_j=x^jyx^{-j}$ to $y_{-j}^{-1}$,
up to the harmless conjugation by a power of $x^n$ in the Schreier
basis. It sends $\Psi_n$ to $-\Psi_n$; the half modulo two of an
even $\Psi_n$ is therefore unchanged.

We also give the direct justification for KM.3's peripheral identity.
Put $g=f(z,x)$ and $h=f(y,z)$, where $z=(xy)^{-1}$. The
multiplier-one three-cycle is $ghf=1$. Right composition with the
actual gentle element $c=(-1,1)$ has normalized pair $(-1,f)$.
Its three-cycle, with $m=-1$, is
\[
 gz^{-1}hy^{-1}fx^{-1}=1.
\]
Substitution of $h=g^{-1}f^{-1}$ gives
$gz^{-1}g^{-1}=xf^{-1}yf=\Phi(z^{-1})$. Consequently
\[
 \Phi(z)=f(z,x)z f(z,x)^{-1}.
\]
This deduction uses the group of actual gentle automorphisms and
the two three-cycle equations; it uses no pentagon or finite lift.

\paragraph{MD.2 (Proved): the off-diagonal beta value is the half-Kummer sign.}
For every $\Phi\in\mathcal H$, in
$R_4=(\mathbb Z/4)[\omega]$, $\omega^2+\omega+1=0$, one has
\[
 \boxed{\quad B_f(-1,\omega)=(-1)^{\Psi_3(\Phi)/2}\quad}.
 \tag{MD.2}
\]
The right side means the two-adic half reduced modulo two.
Its integrality was proved in KM.3, not assumed from arithmetic.
Indeed, the two-cycle gives $A_f(U,V)=A_f(V,U)$.
Reduction modulo two identifies the continuous evaluation maps
$U\mapsto-1$ and $U\mapsto1$, even for the full completed
profinite group ring. Therefore
\[
\begin{split}
 B_f(-1,\omega)
 &=1-2(\omega-1)A_f(-1,\omega)\\
 &=1-2(\omega-1)A_f(1,\omega)\pmod4\\
 &=1-2t_f,
 \qquad t_f=(\omega-1)A_f(\omega,1).
\end{split}
\]
KM.3 gives $t_f\in\mathbb Z_2$ and $\Psi_{3,2}/2=t_f/3$.
Since $3\equiv1\pmod2$, the last displayed unit is precisely
the sign in (MD.2). No gamma factorization is needed.

\paragraph{MD.3 (Proved): the order-three part of duplication is already pentagonal.}
Use the affine realization of $S_3$ over $\mathbb F_3$ with
$p=(1,1)$ and $t=(0,-1)$. Let $a_f$ be defined by
$f(p,t)=p^{a_f}$. Then, for every multiplier-one gentle element,
\[
 a_f=\Psi_2(\Phi)\pmod3.
 \tag{MD.3a}
\]
To prove this directly, evaluate the index-two Schreier basis at
$x=t$, $y=p$. The two conjugates of $y$ have translations $1,-1$,
and their coefficients are $m_0,-m_0$. Thus the translation of
$f(t,p)$ is $2\Psi_2$. The two-cycle also gives
$f(t,p)=p^{-a_f}$. Hence $2\Psi_2=-a_f$ in $\mathbb F_3$,
which is (MD.3a).

For $\Phi\in\mathcal H_5$, the explicit genuine five-point
pentagon calculation FP.33 gives $a_f=\rho\pmod3$. Consequently
\[
 \rho=\Psi_2\pmod3\qquad(\Phi\in\mathcal H_5).
 \tag{MD.3b}
\]
This is only a congruence. It does not identify the entire
frame-defined integer $\rho$ with $\Psi_2$.

\paragraph{MD.4 (Proved reduction; the vanishing assertions remain open).}
For $\Phi\in\mathcal H_5$, define the mixed duplication unit
\[
 D_2(\Phi)=
 \omega^{2\Psi_2(\Phi)}
       \frac{B_f(\omega,\omega)}{B_f(-1,\omega)}\in R_4^\times.
 \tag{MD.4a}
\]
Combining MD.2, MD.3 and the intrinsic phase calculation FP.34 gives
\[
 \boxed{\quad
 D_2(\Phi)=(-1)^{d_2(\Phi)},\qquad
 d_2(\Phi)=\epsilon_2(W_\Phi)+\Psi_3(\Phi)/2\pmod2.
 \quad}
 \tag{MD.4b}
\]
In fact FP.34 says
$B_f(\omega,\omega)=\omega^\rho(-1)^{\epsilon_2(W)}$,
and $2\Psi_2+\rho\equiv0\pmod3$ by MD.3.
Both summands in $d_2$ are continuous additive characters, so
$d_2:\mathcal H_5\to\mathbb F_2$ is an actual continuous
character, independently of any unknown distribution theorem.
The same expression defines a character on all $\mathcal H$.
More precisely, write
$\eta_{\mathrm{aff}}=a_f-\rho=\Psi_2-\rho\pmod3$.
The formula for $D_2$ continues to make sense there and gives
\[
 D_2=\omega^{-\eta_{\mathrm{aff}}}(-1)^{d_2}
 \qquad(\Phi\in\mathcal H).
 \tag{MD.4b$'$}
\]
Indeed $2\Psi_2+\rho=-(\Psi_2-\rho)\pmod3$.
Thus FP.33 is exactly the vanishing of the order-three component
of this specialized duplication unit; its order-two component is
the remaining character $d_2$. In particular $D_2^3=(-1)^{d_2}$
on all $\mathcal H$.

Define a second continuous character on $\mathcal H_5$ by
\[
 e_3(\Phi)=\kappa_d(\Phi)+\Psi_3(\Phi)/2\pmod2.
 \tag{MD.4c}
\]
The five-point frame identity FP.32 then yields the exact decomposition
\[
 \boxed{\quad\Delta|_{\mathcal H_5}=d_2+e_3.\quad}
 \tag{MD.4d}
\]
In particular the main parity equality on $\mathcal H_5$ is
equivalent to $d_2=e_3$ there. Proving both zero would suffice,
but their simultaneous vanishing is stronger than that equivalence.
The equality $d_2=0$ is exactly the single finite specialization
\[
 B_f(\omega,\omega)
       =\omega^{-2\Psi_2(\Phi)}B_f(-1,\omega)\pmod4.
 \tag{MD.4e}
\]
The equality $e_3=0$ is exactly the identification of the
five-point cubic-frame parity with the index-three half-character.
Neither assertion has been deduced here from the full pentagon.
On the arithmetic subgroup $G_{\mathbb Q}\cap\mathcal H$,
Anderson duplication gives $d_2=0$ in the same intrinsic
convention. The supplied arithmetic vanishing of $\Delta$
then gives $e_3=0$ there as well. This uses no identification
of an unnormalized source Kummer exponent with the frame coordinate.

\paragraph{MD.5 (Proved): completed-ring, reflection and cyclotomic audit of KM.5.}
We strengthen the coefficient-level scope of the example, while
keeping it separate from any profinite GT lifting claim.
Write $\widehat{\mathbb Z}=\mathbb Z_2\times H_{\rm odd}$,
where $H_{\rm odd}=\prod_{p\ne2}\mathbb Z_p$.
At a finite odd cyclic quotient, its group algebra over
$\mathbb Z_2$ is finite etale. After the corresponding finite
unramified coefficient extension it splits into its individual
character components. Passing to inverse limits identifies
the full group algebra componentwise with the two-adic
power-series rings in the pro-two exponent variable, with
compatible odd-character labels. An inverse on a character block
is integral whenever its reduction modulo two is a unit.

The idempotents $e_n$ of KM.5 are integral because their
denominators are powers of three. Reduction
$C_{3^{n+1}}\to C_{3^n}$ sends $e_{n+1}$ to $e_n$:
each coset sum acquires the factor three canceled by its
denominator. They therefore define $e$ in
$\mathbb Z_2[[\mathbb Z_3]]$ and, through the direct-factor
inclusion, in $\mathbb Z_2[[\widehat{\mathbb Z}]]$.
On odd-character components $e$ has value one precisely when
the three-primary part of the character has exact order three.
It is zero otherwise, independently of the remaining exponent factors.

For $\Gamma=1-2e$, put
$B(U,V)=\Gamma(U)\Gamma(V)/\Gamma(UV)$.
If the three-primary character of $U$ is trivial then
$\Gamma(U)=1$ and $\Gamma(UV)=\Gamma(V)$, so $B=1$.
The same holds with $U,V$ exchanged. On every remaining
character block both $U-1$ and $V-1$ are integral units:
their nontrivial odd parts have nontrivial reduction modulo two.
Thus there is an integral continuous $A$ with
\[
 B=1+(U-1)(V-1)A,
\]
zero on the former blocks and the indicated quotient on the latter.
Its componentwise definitions are compatible under finite-quotient
transition maps, proving the claimed completed-ring integrality.
They imply $A(U,1)=A(1,V)=0$, and in particular all the
off-diagonal Kummer data of this example vanish.

Every $\lambda\in\widehat{\mathbb Z}^{\times}$ preserves the
exact order of each three-primary character. Therefore
\[
 \Gamma(U^\lambda)=\Gamma(U),\qquad
 B(U^\lambda,V^\lambda)=B(U,V).
\]
The coefficients are in $\mathbb Z_2$ and are unchanged by
unramified coefficient automorphisms. In addition, multiplication
of the exponent by two is an automorphism of the pro-three
factor, so $\Gamma(U^2)=\Gamma(U)$ and
$B(U^2,V^2)=B(U,V)$. Thus, on its odd-root evaluations, the
coefficient-two Frobenius substitution also fixes this example;
the values $\pm1\in\mathbb Z_2$ are fixed by Witt Frobenius.
This is Frobenius substitution, not a multiplication formula
relating shifted arguments. Inversion fixes $e$;
since $\Gamma^2=1$, the reflection identity
$\Gamma(U^{-1})=\Gamma(U)^{-1}$ holds exactly.
Accordingly ordinary cyclotomic substitution covariance and
reflection do not exclude the example. This is a statement about
these specified coefficient identities, not all possible
compatibilities of a GT action on a tower.

The example also obeys both coefficient forms of the
multiplier-one metabelian three-cycle. Indeed $A$ is symmetric,
and, for $W=(UV)^{-1}$, reflection gives
$B(U,V)=B(V,W)=B(W,U)$. Clearing the regular factors
$U-1,V-1,W-1$ gives
\[
\begin{split}
 V^{-1}A(W,U)+(UV)^{-1}A(V,W)+A(U,V)&=0,\\
 V^{-1}A(W,U)+A(V,W)+U A(U,V)&=0.
\end{split}
\]
For example, the first cleared numerator is
$(V-1)/V+(U-1)/(UV)+(W-1)=0$.
The second is $(V-1)/V+(U-1)+U(W-1)=0$.
The relevant factors are regular: on each odd-character block
they are units or nonzero power-series parameters over an
unramified discrete valuation ring.

Nevertheless
\[
 B(\omega,\omega)=-1,\qquad B(-1,\omega)=1,
 \qquad A(\omega,1)=0.
\]
Thus the mixed duplication specialization MD.4e fails when
$\Psi_2=\Psi_3=0$. The exact finite-ring program
\texttt{check\_mixed\_idempotent.py} independently verifies
idempotency, compatibility, integral divisibility, reflection,
all unit substitutions, and the displayed metabelian identities
in $(\mathbb Z/4)[C_{12}]$, $(\mathbb Z/4)[C_{36}]$ and
$(\mathbb Z/4)[C_{108}]$ and their two-variable rings.
It retains the order-four exponent factor through the CRT embedding.
Its output is PASS in each case. The inverse-limit proof above,
not these finitely many checks, proves the completed-ring statement.
No lift of this coefficient datum to a free profinite word satisfying
the full five-point pentagon is asserted.

\paragraph{MD.6 (Known, with its extra hypotheses): the opened duplication theorem.}
Furusho (2003), Theorem 3.1, proves that the relations
$(I),(II),(IV),(K_2)$ imply Anderson duplication $(A_2^0)$;
Theorem 3.4 gives $GT_K\cap I\Gamma\subseteq GT_{A,2^\infty}$.
The explicit group-ring form printed on p.102 is
\[
 B_\sigma\{X,X\}
   =X^{-2\Psi_2^{(0)}(\sigma)}B_\sigma\{-1,X\}.
\]
Here $\Psi_2^{(0)}$ is defined on p.100 by exactly the
index-two Schreier coefficient. In that same coefficient/action
convention, the specialization $X=\omega$ modulo four is
MD.4e. Thus this theorem proves $d_2=0$ for elements satisfying
its extra equations, once those equations are expressed in the
same marking. It supplies neither these extra hypotheses on
all full profinite GT nor the independent comparison $e_3=0$.
An indexed author correction dated 3 January 2008 retains the
same Theorem 3.1 statement; its full PDF could not be retrieved
in this pass, so no detailed claim from the correction is used.

The focused search checked the concrete multiplication and
special-locus leads, including Ihara II, Furusho's author version,
the LNS eigenlocus paper and its conventions appendix, and the
earlier opened NTY comparison. No opened theorem in that search
implies either missing parity comparison on all full profinite GT.
This is a bounded source finding; it is not a theorem of
nonexistence, a proof of independence from the full pentagon,
or a novelty claim.

\paragraph{MD.7 (Proved): two precise equation loci kill the duplication sign.}
Let $C_L$ be the fixed subgroup for the original physical quotient
of A4.M1--A4.M6, rather than the literal current-marking subgroup $C_P$.
Then
\[
 C_L\cap\mathcal H\ \subseteq\ \ker d_2.
 \tag{MD.7a}
\]
Also the multiplier-one locus satisfying the harmonic primitive
equation $(\mathrm{GF}_1)$ lies in $\ker d_2$:
\[
 \{\Phi\in\mathcal H:g=f(x,s)a^{4\rho}\}\subseteq\ker d_2,
 \qquad a=\tau^{-1},\quad g=\beta^{-1}.
 \tag{MD.7b}
\]
Statement A4.M6 identifies this second locus with its precisely marked
NT harmonic equation $(\mathrm I')$.

We prove both inclusions in the projective reduction modulo two
of the affine representation used in FP.34. Put
$c_0=(1-\omega)^{-1}\in\mathbb F_4$. Its relevant values are
\[
 x=(0,\omega),\qquad a=(0,\omega^2),\qquad
 s=(c_0,1),\qquad V=(0,\omega^\rho).
\]
For any $f\in F'$, its evaluation at
$(0,\omega),(c_0,1)$ is the translation $c_0t_f$.
Its evaluation at $(0,\omega^2),(c_0,1)$ is the translation
$c_0\overline{t_f}$. By KM.3, $\overline{t_f}=t_f$ and
$t_f\bmod2=\Psi_3/2\bmod2$.

For $\Phi\in C_L$, equation (M7) in A4.M6 gives
$g=f(a,s)x^{-b_0}$. Since the rightmost factor has zero
translation, the translation of $g$ is $c_0t_f$.
Under $(\mathrm{GF}_1)$ the same conclusion follows from
$g=f(x,s)a^{4\rho}$. Finally $W=gV$, and $V$ also has
zero translation. The proof of FP.34 identifies
$\operatorname{trans}(W)/c_0$ with $\epsilon_2(W)$.
Therefore both hypotheses give
\[
 \epsilon_2(W)=t_f=\Psi_3/2\pmod2,
\]
which proves (MD.7a)--(MD.7b).

Neither argument uses the pentagon, an equality between the
existential GF0 parameter and $\Psi_3$, nor a full Anderson
multiplication law. In particular this proves only the harmonic
half-character comparison on $C_L$. On
$C_L\cap\mathcal H_5$ the remaining main obstruction is exactly
$e_3=\kappa_d+\Psi_3/2$; its vanishing is still open here.

\paragraph{MD.8 (Proved): the half-character is a descended harmonic sign.}
The half-character has a universal interpretation, stronger than
the fixed-locus statement MD.7:
\[
 \boxed{\quad
 \Psi_3(\Phi)/2\bmod2
 =\epsilon_2(W_{D_L\Phi})
 =\epsilon_2(W_{D_P\Phi})
 \qquad(\Phi\in\mathcal H).
 \quad}
 \tag{MD.8a}
\]
Consequently, as characters on $\mathcal H$,
\[
 d_2=\epsilon_2\circ W+\epsilon_2\circ W\circ D_L
     =\epsilon_2\circ W+\epsilon_2\circ W\circ D_P.
 \tag{MD.8b}
\]
In particular both $C_L\cap\mathcal H$ and
$C_P\cap\mathcal H$ lie in $\ker d_2$. The two fixed groups
themselves remain distinct; this statement does not identify them.

For the proof, keep the same affine quotient and $c_0$ as in MD.7.
The established pure normalizations of the two descents are
\[
 \beta_{D_L}=x^{b_0}P^{-1},\quad P=f(a,s),\qquad
 \beta_{D_P}=x^{a_0}C,\quad C=f(s,\tau).
\]
The image of $x$ has zero translation, while
\[
 \operatorname{trans}(P)=c_0t_f,\qquad
 \operatorname{trans}(C)=-c_0t_f.
\]
The first identity was proved in MD.7. For the second, the
modulo-two images are $s=(c_0,1)$ and $\tau=(0,\omega)$;
the two-cycle turns the evaluation into the inverse of the
one defining $c_0t_f$.
Taking inverses of the displayed pure normalizations, the
rightmost powers of $x$ still have zero translation. It follows
that both $\beta_{D_L}^{-1}$ and $\beta_{D_P}^{-1}$ have
translation $c_0t_f$. The respective normalizing shears
$V_{D_L\Phi}$ and $V_{D_P\Phi}$ also have zero translation.
Thus both corrected harmonic frames have translation $c_0t_f$,
and FP.34 gives (MD.8a). No comparison of the integer gauges
$a_0,b_0$ or of the individual old/new Galois labels is required.

\paragraph{MD.9 (Proved): the square-pair comparison is not a word identity.}
The following exact finite lemma isolates an additional issue in
the modulo-four Burau pentagon comparison. Let
\[
 p=(123),\qquad t=(1234)\quad\text{in }S_4.
\]
Then $(p,t)$ has peripheral orders $(3,4,4)$ and generates $S_4$,
whereas $(p,t^2)$ has peripheral orders $(3,2,3)$ and generates
$A_4$. The joint same-word evaluation
\[
 F\longrightarrow S_4\times A_4,
 \qquad x\longmapsto(p,p),\quad y\longmapsto(t,t^2)
 \tag{MD.9a}
\]
is onto. Consequently its restrictions have images
\[
 F'\longmapsto A_4\times V_4,
 \qquad F''\longmapsto V_4\times\{1\}.
 \tag{MD.9b}
\]

For the proof, both projections in (MD.9a) are onto. The nontrivial
quotients of $S_4$ are $S_4,S_3,C_2$, and those of $A_4$ are
$A_4,C_3$. The groups have no common nontrivial quotient.
Goursat's lemma therefore gives the full direct product.
Taking its first and second derived subgroups proves (MD.9b).
The discrete word
\[
 w=[[x,y],x[x,y]x^{-1}]
\]
provides a concrete check: $w(p,t)=(14)(23)$ but
$w(p,t^2)=1$. The independent exact permutation program
\texttt{check\_square\_pair\_goursat.py} verifies the pair orders,
the joint derived-series orders $288,48,4,1$, and this witness.

The square pair is relevant to the preceding finite Burau
calculation for a precise reason. After removing its central
order-three factor, that calculation gives a group generated by
$p,t$ with peripheral orders $(3,4,4)$. Its square-pair subgroup
is $A_4\times C_2$; if $e$ denotes its central involution, then
$p,et^2$ generate an $A_4$ complement with orders $(3,2,3)$.
For every $f\in F'$, the central factor cancels and
$f(p,t^2)=f(p,et^2)$. This identifies the square-pair value with
the half-Kummer $A_4$ evaluation, after the stated marked
isomorphism. The central decomposition and that marked
isomorphism are finite-Burau inputs, not consequences of
Goursat's lemma.

Thus an equality comparing the relevant coordinate of $f(p,t)$
with the half-Kummer value of $f(p,t^2)$ is a substantive
distribution statement; it cannot be inferred merely from
using a single commutator word in both places. The lemma does
not assert that $w$ satisfies the gentle equations, nor that
either independent finite assignment lifts to a genuine GT
element. It therefore does not prove independence from the
gentle equations or the full profinite pentagon.

\paragraph{MD.10 (Proved using LI.8): duplication alone cannot settle the gentle cubic sign.}
Extend the second defect to the actual gentle group by
\[
 \widetilde e_3(\Phi)=\epsilon_3(T_\Phi)+\Psi_3(\Phi)/2\pmod2.
\]
Its restriction to $\mathcal H_5$ is $e_3$ by FP.32.
The actual gentle element $k$ supplied by LI.8 satisfies
\[
 B_k=1,\qquad \Psi_n(k)=0\ (n\geq1),\qquad
 \rho(k)=0,\qquad d_2(k)=0,\qquad
 \widetilde e_3(k)=\Delta(k)=1.
 \tag{MD.10}
\]
Indeed LI.8 gives $f_k,\beta_k\in R_{\rm nil}(F'')\subseteq F'''$
and $\Delta(k)=1$. In particular $f_k$ has zero image in
$F'/F''$, proving $A_{f_k}=0$ and $B_k=1$.
For every $n$, the Schreier subgroup $N_n$ contains $F'$, so
$F''\subseteq N_n'$ and every coefficient defining $\Psi_n(k)$
vanishes. The abelianization of $\beta_k$ is zero, giving
$\rho(k)=0$. Formula FP.34 now gives $\epsilon_2(W_k)=0$,
and the other assertions follow.

Consequently this element satisfies the beta-factorization,
reflection and multiplication coefficient equations with the
choice $\Gamma_k=1$ and all Kummer exponents zero, while retaining
the nontrivial cubic defect. This does not put $k$ in $GT_A$,
whose definition already requires membership in full profinite GT.
The full pentagon for $k$ remains unproved. The conclusion is
instead the precise logical limitation that a beta/Kummer
comparison, including the complete vanishing of $d_2$, still
needs a cubic-frame comparison to establish the desired parity.
